\documentclass[12pt]{article}
\usepackage[margin=1in]{geometry}
\usepackage[backend=biber, style=apa, sorting=nyt]{biblatex}
\usepackage{hyperref}
\usepackage{parskip}
\usepackage{setspace}
\addbibresource{references.bib}

\title{\textbf{Annotated Bibliography} \\
\large Predicting Whole-Brain Neural Dynamics from Task-Based fNIRS Signals: \\
Attentional Stability as a Bridge Between Structured Tasks and Naturalistic Brain Activity}
\author{MACS 30200: Computational Research Design}
\date{Spring 2026}

\begin{document}
\maketitle
\onehalfspacing

\section*{Overview}

This bibliography is organized around a central research question: does the consistency with which a person pays attention during simple lab tasks tell us anything about how well a computational model can reconstruct that person's brain activity during something far more complex, like watching a movie?

To understand why this question matters, it helps to know the difference between the two brain imaging tools at the center of this project. Functional near-infrared spectroscopy (fNIRS) is a lightweight, wearable headset that measures brain activity by detecting how much infrared light is absorbed through the scalp. It is relatively affordable and can be used outside a lab, but it only captures signals from the surface of the brain. Functional magnetic resonance imaging (fMRI), on the other hand, uses a large scanner to measure activity across the whole brain, including deep structures, but it is expensive, loud, and requires participants to lie completely still. The core appeal of this project is that if we can build a model to translate portable fNIRS signals into fMRI-quality brain maps, it would open up whole-brain neuroscience to settings where fMRI is simply not practical.

This project builds on Gao et al. (2025), who showed that such a translation is possible when both imaging types are recorded during the same movie. The proposed work extends this by asking two harder questions. First, can the translation still work when the fNIRS signals come from structured cognitive tasks, like a memory test or an attention task, rather than from movie-watching? Second, do people who respond more consistently during those tasks also produce fNIRS signals that translate more accurately? That consistency is measured through reaction time variability, which captures how evenly paced a person's button-press responses are across many trials. Someone with low variability responds with roughly the same speed each time; someone with high variability has a mix of fast and slow responses, suggesting their attention is drifting in and out. The sources in this bibliography are grouped into three areas: the technical feasibility of fNIRS-to-fMRI translation, the science of reaction time variability as a measure of attentional control, and what we know about how the brain flexibly adapts its activity across different situations.

\textbf{Significance of the Study:}
Right now, studying the brain in real-world social contexts: classrooms, workplaces, clinical settings — is practically impossible with fMRI because it requires a multi-million dollar machine and a participant who can lie completely still. This means most of what we know about the neural basis of social behavior, attention, and mental health comes from highly artificial lab conditions that may not reflect how people actually think and feel in their lives. fNIRS is portable and affordable, but it only sees the surface of the brain, which limits how much you can actually learn from it.

If this project works, a researcher could deploy a cheap fNIRS headset in a real classroom or clinic, collect brain data during naturalistic tasks, and use the model to reconstruct what the whole brain was likely doing: fMRI-level insight without fMRI infrastructure. If cross-task translation works well, it validates structured lab tasks as affordable proxies for capturing whole-brain dynamics, with direct implications for mental health screening and educational neuroscience in under-resourced settings. If attentional stability predicts translation accuracy, it adds another layer: the model not only reconstructs brain activity, it also reveals something meaningful about the individual's cognitive profile from behavioral data alone.

\bigskip

%% -------------------------------------------------------
%% PILLAR 1: fNIRS-fMRI NEURAL TRANSLATION
%% -------------------------------------------------------
\section*{Pillar 1: fNIRS-to-fMRI Neural Translation}

\subsection*{\cite{gao2025}}
\textbf{Summary:} Gao et al. build a computational model that takes fNIRS signals from the prefrontal cortex (the front region of the brain, involved in attention, planning, and social reasoning) and uses them to predict activity across the whole brain as measured by fMRI, all while participants watch a TV episode. The model first reduces both the fNIRS and fMRI recordings to their most statistically important underlying components using principal component analysis, then learns a linear mapping between those components. Training was done on the first half of a BBC \textit{Sherlock} episode, and testing was done on a held-out participant watching the second half, so neither the person nor the content had been seen by the model before. The model successfully predicted fMRI activity in 66 of 122 brain regions, including deep subcortical structures like the basal ganglia and hippocampus that fNIRS physically cannot reach. Performance was strongest in the default mode network, which is active during internally directed thought and narrative processing, and in the frontoparietal control network, involved in flexible goal-directed cognition. The model also retained meaningful information about the semantic content of the movie, not just the timing of activity, and generalized to a different TV show (\textit{Friday Night Lights}), suggesting the learned mapping is not specific to one stimulus.

\textbf{Evaluation:} This paper is the direct predecessor to this project and is probably the most important source in this bibliography. Its central contribution is demonstrating that cross-participant, cross-stimulus fNIRS-to-fMRI translation is computationally achievable even when the two participant groups were never scanned together. This project makes the problem harder by using fNIRS recorded during structured cognitive tasks rather than during movie-watching, dropping the assumption that both modalities share the same stimulus. One notable gap in Gao et al. is that they do not investigate why the model works better for some participants than others. This project takes that gap as its starting point by asking whether individual differences in attentional stability, measured through reaction time variability, explain variation in translation accuracy.

\subsection*{\cite{burns2019}}
\textbf{Summary:} Burns et al. used fNIRS to measure prefrontal brain responses to persuasive health messages in a sample from the Middle East, a population that is heavily underrepresented in neuroscience research. The study showed that fNIRS can detect meaningful differences in neural responses to persuasive content in a real-world, cross-cultural setting outside the laboratory, and that those prefrontal signals predicted participants' behavioral intentions toward the health recommendations. The focus on prefrontal cortex coverage is directly relevant because this project also relies on prefrontal fNIRS.

\textbf{Evaluation:} This paper is included primarily as a methodological anchor. It demonstrates that fNIRS recordings from the prefrontal cortex are sensitive enough to detect meaningful individual differences in brain responses to socially and cognitively demanding content, which is a core assumption this project depends on. It also involves Shannon Burns, one of the co-authors of Gao et al. (2025), and uses comparable fNIRS hardware and data collection protocols, making it a useful reference point for understanding the quality and scope of data produced by this research group. More broadly, it supports the argument that fNIRS is a credible tool for social and cognitive neuroscience beyond tightly controlled laboratory paradigms.

\subsection*{\cite{liu2015}}
\textbf{Summary:} Liu et al. asked whether fNIRS signals recorded from the scalp could be used to infer fMRI activity in deep brain structures that fNIRS physically cannot reach. They scanned the same participants simultaneously with both devices across several cognitive tasks, including a go/no-go response inhibition task and an emotional face recognition task, and found that activity in deep regions including the insula (involved in bodily awareness and emotional processing), the hippocampus (involved in memory), and the caudate (a structure involved in habit and reward processing) could be predicted from fNIRS signals measured at the cortical surface. The prediction was possible because these deep regions are functionally connected to the prefrontal surface regions that fNIRS does capture, meaning their activity tends to go up and down together.

\textbf{Evaluation:} This is an important precursor to Gao et al. (2025) and to this project because it establishes the neurobiological basis for why fNIRS-to-fMRI inference works at all. If deep and surface brain regions were not functionally linked, there would be no information in the fNIRS signal to exploit. The key limitation Liu et al. acknowledge is that their model required simultaneous fNIRS and fMRI recording from the same participants, which is expensive and practically difficult. Gao et al. (2025) solved this by using a shared movie stimulus to align signals across separately scanned groups. This project inherits that solution and extends it further by asking whether the alignment can cross task types, not just stimulus instances.

%% -------------------------------------------------------
%% PILLAR 2: RT VARIABILITY AND ATTENTIONAL STABILITY
%% -------------------------------------------------------
\section*{Pillar 2: Reaction Time Variability and Attentional Stability}

\subsection*{\cite{kofler2013}}
\textbf{Summary:} This is a meta-analysis, meaning it statistically combines findings across many separate studies rather than reporting new data. Kofler et al. pooled results from 319 studies on reaction time variability in people with ADHD (Attention Deficit Hyperactivity Disorder). Reaction time variability refers to how inconsistently a person responds across trials of a task: someone with low variability presses a button at roughly the same speed each time, while someone with high variability produces a mix of fast and slow responses, suggesting their attention is drifting. The review found large and reliable differences between people with ADHD and typically developing individuals (effect size Hedges' \textit{g} = 0.76 for children, 0.46 for adults), and importantly, reaction time variability was a stronger discriminator than average response speed. Stimulant medications reduced variability substantially (g = 0.74), but behavioral interventions did not. Elevated variability appeared consistently across many task types including the N-back memory task and the Continuous Performance Task, both of which are present in the dataset this project uses.

\textbf{Evaluation:} This is the primary source justifying the choice of reaction time variability as the individual difference measure in this project. Three things stand out. First, variability is more theoretically meaningful than mean response speed because it reflects moment-to-moment fluctuations in attentional control, which is exactly what this project wants to measure. Second, the fact that variability appears on both the N-back and Continuous Performance Task means it can be computed from the data already available. Third, elevated variability is not unique to ADHD but appears across a range of clinical and subclinical populations, which supports treating it as a continuous dimension across all individuals rather than a binary clinical label. This paper also satisfies the course requirement for at least one review article.

\subsection*{\cite{sonugabarke2007}}
\textbf{Summary:} Sonuga-Barke and Castellanos propose what they call the Default Mode Interference hypothesis to explain why reaction time variability occurs neurologically. The argument starts from a well-established observation: the brain has a ``default mode network,'' a set of interconnected regions that become active when a person is resting, daydreaming, or thinking about other people. During a goal-directed task, this network is normally suppressed so that task-relevant regions can operate without interference. The paper argues that in people with attentional difficulties, this suppression fails periodically. When the default mode network briefly reactivates during a task, it disrupts the processing of the current stimulus, causing the person's attention to lapse for a moment. That lapse shows up behaviorally as an unusually slow response on that particular trial, producing the pattern of occasional very long reaction times that defines high reaction time variability.

\textbf{Evaluation:} This paper is probably the most theoretically important source in the bibliography because it connects the behavioral measure this project uses (reaction time variability) to the neural dynamics the model is trying to predict. The default mode network is the same network where Gao et al. (2025) found the strongest fNIRS-to-fMRI translation accuracy, and it is also the network that is especially active during naturalistic experiences like watching a movie. If people with high reaction time variability have poor default mode suppression during their cognitive tasks, then their task-based fNIRS signals will contain more ``rest-like'' activity mixed into what should be ``task-like'' activity, and the model should find it harder to translate those signals into the correct movie-watching brain response. This gives the project a specific, testable neural mechanism, not just a behavioral correlation.

\subsection*{\cite{macdonald2009}}
\textbf{Summary:} MacDonald, Li, and Backman review a body of research on within-person variability in cognitive performance and argue that how consistently a person performs across trials reflects the health and efficiency of their underlying neural systems. They identify three biological mechanisms that drive this consistency: the stability of dopamine signaling (dopamine is a chemical messenger in the brain that affects how reliably neurons fire), the integrity of white matter (the long-range connections that link different brain regions to each other), and the regulatory strength of the prefrontal cortex. When any of these systems is less efficient, the brain produces noisier neural signals, which in turn produces more erratic behavioral outputs.

\textbf{Evaluation:} This review is useful for grounding the choice of reaction time variability as an individual difference measure in something more substantive than a statistical observation. For this project, the argument that RT variability tracks individual differences in neural signal quality is directly relevant: a person whose brain generates cleaner, more stable neural signals should also produce more consistent fNIRS recordings, which in turn should be easier for a computational model to translate into an accurate fMRI-like output. The paper also reinforces the case for treating RT variability as a continuous measure that varies meaningfully across all healthy people, not just those with a clinical diagnosis. This matters because the dataset for this project consists of healthy participants with no clinical information.

%% -------------------------------------------------------
%% PILLAR 3: NEURAL FLEXIBILITY AND CROSS-CONTEXT DYNAMICS
%% -------------------------------------------------------
\section*{Pillar 3: Neural Network Flexibility and Cross-Context Dynamics}

\subsection*{\cite{cole2013}}
\textbf{Summary:} Cole et al. examined how the brain reorganizes its communication patterns when people switch between 64 different cognitive tasks. Brain connectivity, in this context, means how strongly two regions co-activate or coordinate with each other during a given moment. They found that a set of regions in the frontoparietal network, a distributed network spanning the front and side of the brain that supports flexible, goal-directed thinking, acted as ``flexible hubs.'' These hubs shifted their connectivity patterns most dramatically across task conditions, while regions in other networks (such as the visual or somatomotor networks) remained relatively stable regardless of what task was being performed. The connectivity patterns of these hubs could reliably identify which task a person was currently doing, and this held even for tasks the person had never practiced before.

\textbf{Evaluation:} This paper provides the neural architecture argument underpinning one of the key predictions in this project. The frontoparietal network substantially overlaps with the prefrontal cortex, which is the exact region fNIRS measures. If some people's prefrontal fNIRS signals more fully reflect the flexible hub reorganization described here, those signals should carry richer, more generalizable information about the brain's current cognitive state, making them easier to translate into the movie-watching fMRI response. The paper also sets up a specific prediction about where the model should perform best and worst. Strong performance in the frontoparietal and default mode networks, and weak performance in the somatomotor network (which handles physical movement and sensation), would be consistent with both Cole et al. and Gao et al. (2025).

\subsection*{\cite{braun2015}}
\textbf{Summary:} Braun et al. scanned 344 healthy adults with fMRI while they performed N-back tasks at different difficulty levels. The N-back task asks participants to watch a sequence of items appear one at a time and indicate whether each item matches the one that appeared a certain number of steps earlier. In the 2-back version, participants must hold two prior items in mind simultaneously while processing the current one. The study found that how much the frontal brain networks reorganized their communication patterns during the 2-back condition predicted two outcomes: accuracy on the task itself, and performance on the Trail Making Test B, a standard out-of-scanner measure of cognitive flexibility that requires alternating between numbers and letters. Participants whose frontal networks showed greater cross-regional integration during the demanding N-back condition performed better on both.

\textbf{Evaluation:} This paper is relevant on two levels. First, it validates the N-back task as a meaningful window into cognitive flexibility more broadly, not just working memory capacity, which supports using N-back reaction times as the basis for the attentional stability measure in this project. Second, it demonstrates empirically that individual differences in how the brain organizes itself during a structured task carry information about cognitive functioning outside that task. This is the same logic this project applies: people whose brains are more flexibly organized during the N-back task may also produce fNIRS signals that translate more accurately into the movie-watching brain response.

\subsection*{\cite{tavor2016}}
\textbf{Summary:} Tavor et al. published this paper in \textit{Science} and asked a question that is closely related to what this project is trying to do: can brain activity measured in one context predict brain activity in a completely different context? Specifically, they used resting-state fMRI (scans taken while participants simply lie still and let their mind wander) from 98 participants drawn from the Human Connectome Project, a large publicly available brain imaging dataset, to train a model that predicted how each individual's brain would respond during seven different cognitive tasks. The model was linear and required no task data from the individuals being predicted. It worked well across most tasks and most brain regions, suggesting that the brain's baseline connectivity structure, captured at rest, encodes a kind of blueprint for how it will respond to external demands.

\textbf{Evaluation:} Out of all the sources here, this paper sits closest to what this project is trying to do. Tavor et al. predict across contexts using learned computational mappings, and so does this project, though in a different direction: from structured tasks to movie-watching rather than from rest to tasks. What is especially useful is that Tavor et al. note meaningful individual differences in prediction accuracy across participants, with some people being consistently easier to predict than others. That observation is the empirical starting point for the central question in this project: whether a person's attentional stability, measured by reaction time variability, helps explain how predictable their brain is across different cognitive contexts.

\printbibliography[title={References}]

\end{document}
